\subsection{伽罗瓦代数}

设 K 是一个交换域。若 E 是一个 K-代数，则我们用 \(\operatorname{Aut}_{\mathrm{K}}(\mathrm{E})\) 表示它的自同构群。若 E 是体 K 的一个伽罗瓦扩张，则群 \(\operatorname{Aut}_{\mathrm{K}}(\mathrm{E})\) 就是伽罗瓦群 \(\operatorname{Gal}(\mathrm{E}/\mathrm{K})\)（V, §10, No.~2, 第 58 页）。

设 G 是一个群。一个 \((K,G)\)-代数是装备了一个群同态 \(\lambda:G\to\operatorname{Aut}_{K}(E)\) 的 K-代数 E。于是同态 \(\lambda\) 给 E 装备了一个以 G 为算子域的带算子群的结构，以及一个左 K{[}G{]}-模的结构，其外部律由下式给出

\[
\Bigl (\sum_ {g \in \mathrm{G}} \mu_ {g} g \Bigr) x = \sum_ {g \in \mathrm{G}} \mu_ {g} \lambda (g). x\tag{14}
\]

对每个 \(x \in \mathrm{L}\) 与每个元素 \((\mu_g)_{g \in \mathrm{G}}\)（\(\mathrm{K}[\mathrm{G}]\) 的）。一个 \((\mathrm{K}, \mathrm{G})\)-代数的态射是一个代数态射，它同时也是一个带算子群的态射。

对每个族 \((\mathrm{E}_{i})_{i\in\mathrm{I}}\)（\((\mathrm{K},\mathrm{G})\)-代数的），装备其带算子群的积群结构的积 K-代数 \(\prod_{i\in I}E_{i}\) 是一个 \((\mathrm{K},\mathrm{G})\)-代数。

若 \(\mathrm{E}\) 是一个 (K, G)-代数，则集合 \(\mathrm{E}^{\mathrm{G}}\)（由 \(\mathrm{E}\) 的在 \(\mathrm{G}\) 下不变的元素组成）是 \(\mathrm{E}\) 的一个子代数。

设 E 是一个 (K, G)-代数，其中 G 经 \(\lambda: G \to \operatorname{Aut}_{\mathrm{K}}(\mathrm{E})\) 作用。若 \(K'\) 是 K 的一个扩张，则对任何 \(g \in G\)，用 \(\lambda'(g)\) 表示自同构 \(\operatorname{Id}_{\mathrm{K}'} \otimes \lambda(g)\)（\(K'\)-代数 \(\operatorname{L}_{(\mathrm{K}')}\) 的）。于是 \(\lambda': g \mapsto \lambda'(g)\) 给 \(\operatorname{L}_{(\mathrm{K}')}\) 装备了一个 \((\mathrm{K}', \mathrm{G})\)-代数的结构。

给定一个群 H 以及 H-集 X 与 Y，我们用 \(\mathcal{F}_{\mathrm{H}}(\mathrm{X},\mathrm{Y})\) 表示从 X 到 Y 的 H-集态射的集合。它是映射 \(f:X\to Y\)（满足 \(f(hx)=hf(x)\)，对每个 \(h\in H\) 与 \(x\in X\)）的集合。

设 G 是一个有单位元素 e 的群，设 H 是 G 的一个子群，并设 E 是一个 (K,H)-代数。由 E 经从 H 到 G 的余诱导得到的 (K,G)-代数，记作 \(\mathrm{Coind}_{\mathrm{H}}^{\mathrm{G}}(\mathrm{E})\)，是装备了 G 的一个作用的 K-代数 \(\mathcal{F}_{\mathrm{H}}(\mathrm{G},\mathrm{E})\)，该作用由同态 \(\lambda\)（从 G 到 \(\mathrm{Aut}_{\mathrm{K}}(\mathrm{Coind}_{\mathrm{H}}^{\mathrm{G}}(\mathrm{E}))\)，由下式定义）给出

\[
(\lambda (g) \cdot f) (g ^ {\prime}) = f (g ^ {\prime} g)\tag{15}
\]

对 \(f \in \text{Coind}_{\text{H}}^{\text{G}}(\text{E})\) 与 \(g, g' \in G\)。

引理 7. —— a) 设 S 是 G 的一个子集，使得 G 的每个元素都可以唯一地写成 hs（其中 \(h \in H\) 与 \(s \in S\)）。于是把 f 映到 \((f(s))_{s \in S}\) 的映射是从 K-代数 \(\operatorname{Coind}_{\mathrm{H}}^{\mathrm{G}}(\mathrm{E})\) 到由 S 到 E 的映射组成的 K-代数 \(\mathcal{F}(\mathrm{S}, \mathrm{E})\) 的一个同构。

\begin{enumerate}
\def\labelenumi{\alph{enumi})}
\setcounter{enumi}{1}
\tightlist
\item
  代数 \(\operatorname{Coind}_{\mathrm{H}}^{\mathrm{G}}(\mathrm{E})\) 在 K 上有有限次数，当且仅当 E 在 K 上有有限次数且 H 在 G 中的指数有限。此时，我们有公式
\end{enumerate}

\[
[ \mathrm{Coind} _ {\mathrm{H}} ^ {\mathrm{G}} (\mathrm{E}): \mathrm{K} ] = (\mathrm{G}: \mathrm{H}) [ \mathrm{E}: \mathrm{K} ].
\]

\begin{enumerate}
\def\labelenumi{\alph{enumi})}
\setcounter{enumi}{2}
\tightlist
\item
  设 \(E^{H}\) 是群 H 在 E 中的不变量代数，而 \(\operatorname{Coind}_{\mathrm{H}}^{\mathrm{G}}(\mathrm{E})^{\mathrm{G}}\) 是群 G 在 \(\operatorname{Coind}_{\mathrm{H}}^{\mathrm{G}}(\mathrm{E})\) 中的不变量代数。映射 \(f \mapsto f(e)\)
\end{enumerate}

从 \(\mathrm{Coind}_{\mathrm{H}}^{\mathrm{G}}(\mathrm{E})\) 到 E 限制为从 \(\mathrm{Coind}_{\mathrm{H}}^{\mathrm{G}}(\mathrm{E})^{\mathrm{G}}\) 到 \(\mathrm{E}^{\mathrm{H}}\) 的一个代数同构。

断言 a) 由定义得出并蕴含 b)。由 (15)，一个从 G 到 E 的映射是 \(\mathrm{Coind}_{\mathrm{H}}^{\mathrm{G}}(\mathrm{E})\) 的在 G 下不变的元素，当且仅当它是常值等于 E 的一个在 H 下不变的元素的映射。

注 1. —— 设 G 是一个群，H 是 G 的一个子群，而 N 是 H 的一个子群。设 E 是一个 (K,N)-代数。设 \(\alpha\) 是 \(\mathcal{F}_{\mathrm{H}}(\mathrm{G},\mathcal{F}_{\mathrm{N}}(\mathrm{H},\mathrm{E}))\) 的一个元素。我们有关系

\[
\alpha (g) (n h) = n (\alpha (g) (h))\tag{16}
\]

对每个 \(g \in \mathbf{G}\)、\(h \in \mathbf{H}\) 与 \(n \in \mathbf{N}\)，以及关系

\[
\alpha (h g) (h ^ {\prime}) = (h \alpha (g)) (h ^ {\prime}) = \alpha (g) (h ^ {\prime} h)\tag{17}
\]

对所有 \(h, h' \in H\) 与每个 \(g \in G\)。因此，\(\alpha(ng)(e) = \alpha(g)(n) = n(\alpha(g)(e))\)（对每个 \(g \in G\) 与 \(n \in N\)）。于是我们可以考虑映射

\[
\psi : \mathcal {F} _ {\mathrm{H}} (\mathrm{G}, \mathcal {F} _ {\mathrm{N}} (\mathrm{H}, \mathrm{E})) \to \mathcal {F} _ {\mathrm{N}} (\mathrm{G}, \mathrm{E})
\]

由关系式 \(\psi(\alpha)(g)=\alpha(g)(e)\)（其中 \(\alpha\) 取遍 \(\mathcal{F}_{\mathrm{H}}(\mathrm{G},\mathcal{F}_{\mathrm{N}}(\mathrm{H},\mathrm{E}))\)，g 取遍 G）定义。映射 \(\psi\) 是从 \(\operatorname{Coind}_{\mathrm{H}}^{\mathrm{G}}(\operatorname{Coind}_{\mathrm{N}}^{\mathrm{H}}(\operatorname{E}))\) 到 \(\operatorname{Coind}_{\mathrm{N}}^{\mathrm{G}}(\operatorname{E})\) 的一个代数同构，其逆把元素 \(\beta\)（\(\mathcal{F}_{\mathrm{N}}(\mathrm{G},\mathrm{E})\) 的）映到映射 \(\alpha\)，后者由关系式 \(\alpha(g)(h)=\beta(hg)\)（对 \(g\in G\) 与 \(h\in H\)）定义。

我们现在假设给定一个有限群 G 以及一个约化的（V, §6, No.~6, 第 32 页）有限次交换 K-代数 L，它装备了由同态 \(\lambda\)（从 G 到 \(\mathrm{Aut}_{\mathrm{K}}(\mathrm{L})\)）所给出的 G 的一个作用。对 \(x \in \mathrm{L}\) 与 \(g \in \mathrm{G}\)，我们用 \(g \cdot x\) 表示 \(x\) 在 L 的自同构 \(\lambda(g)\) 下的变换。设 \(\mathcal{S}\) 是 L 的极大理想的集合；我们用 \(g \cdot \mathfrak{m}\) 表示元素 \(\mathfrak{m}\)（\(\mathcal{S}\) 的）在 L 的自同构 \(\lambda(g)\) 下的变换。它是 \(\mathcal{S}\) 的一个元素。对每个 \(\mathfrak{m}\)（\(\mathcal{S}\) 中的），体 L/ \(\mathfrak{m}\) 是 K 的一个有限扩张。我们用 \(\pi_{\mathfrak{m}}: \mathrm{L} \to \mathrm{L}/\mathfrak{m}\) 表示投影，并用 \(\mathrm{G}_{\mathfrak{m}}\) 表示 \(\mathfrak{m}\) 在 G 中的稳定子，即这样的 \(g \in \mathrm{G}\)（使得 \(g \cdot \mathfrak{m} = \mathfrak{m}\)）组成的集合。K-代数 L/ \(\mathfrak{m}\) 装备了 \(\mathrm{G}_{\mathfrak{m}}\) 的一个作用，它经由同态 \(\lambda_{\mathfrak{m}}\)（从 \(\mathrm{G}_{\mathfrak{m}}\) 到 \(\mathrm{Aut}_{\mathrm{K}}(\mathrm{L}/\mathfrak{m})\)）给出，该同态把元素 \(h\)（\(\mathrm{G}_{\mathfrak{m}}\) 的）映到 L/ \(\mathfrak{m}\) 的、由 \(\lambda(h)\) 通过过渡到商所导出的自同构。

设 \(\mathcal{O}\) 是 G 在 \(\mathcal{S}\) 中的轨道的集合。给定一条轨道 \(\sigma \in \mathcal{O}\)，令 \(\mathfrak{a}_{\sigma} = \bigcap_{\mathfrak{m} \in \sigma} \mathfrak{m}\) 与 \(\mathrm{L}_{\sigma} = \mathrm{L} / \mathfrak{a}_{\sigma}\)。由于 \(\mathfrak{a}_{\sigma}\) 在 G 下不变，通过过渡到商，G 在 L 上的作用定义了一个同态 \(\lambda_{\sigma}\)（从 G 到 \(\mathrm{Aut}_{\mathrm{K}}(\mathrm{L}_{\sigma})\)）。最后，用 \(\pi_{\sigma}\) 表示从 L 到 \(\mathrm{L}_{\sigma}\) 的典范映射。

引理 8. —— a) 对每个 \(g \in G\)、\(\sigma \in \mathcal{O}\) 与 \(\mathfrak{m} \in \sigma\)，我们有

\[
[ \mathrm{L} _ {\sigma}: \mathrm{K} ] = \operatorname{Card} (\sigma) [ \mathrm{L} / \mathfrak {m}: \mathrm{K} ].
\]

\begin{enumerate}
\def\labelenumi{\alph{enumi})}
\setcounter{enumi}{1}
\item
  映射 \(\pi : x \mapsto (\pi_{\sigma}(x))_{\sigma \in \mathcal{O}}\) 是从 L 到 \(\prod_{\sigma \in \mathcal{O}} \mathrm{L}_{\sigma}\) 的一个 (K,G)-代数同构。
\item
  用 \(L^{G}\)（相应地，\(L_{\sigma}^{G}\)）表示 L（相应地，\(L_{\sigma}\)）的在 G 的作用下不变的元素组成的子代数。于是 \(\pi\) 诱导一个从 \(L^{G}\) 到 \(\prod_{\sigma\in\mathcal{O}}L_{\sigma}^{G}\) 的同构。
\end{enumerate}

由于代数 L 是约化的且有限次，L 的极大理想的交化归为 0（VIII, 第 173 页, 推论 2）。此外，若 m 与 \(m'\) 是 L 的两个不同的极大理想，则我们有 \(m + m' = L\)。由 I, §8, No.~11, 第 110 页的命题 10，从 L 到 \(\prod_{m \in S} L/m\) 的典范映射是一个同构，且从 \(L/a_{\sigma}\) 到 \(\prod_{m \in \sigma} L/m\) 的典范映射（对每个 \(\sigma \in O\)）也是同构。断言 a) 由此得出。由于 O 是 S 的一个分拆，断言 b) 成立；断言 c) 是 b) 的直接推论。

我们现在固定一条轨道 \(\sigma \in \mathcal{O}\) 与一个元素 \(\mathfrak{m}\)（\(\sigma\) 的）。令 \(F_{\mathfrak{m}} = \text{Coind}_{G_{\mathfrak{m}}}^{G}(L/\mathfrak{m})\)，并用 \(\lambda_{F_{\mathfrak{m}}}\) 表示 \(G\) 在 \(F_{\mathfrak{m}}\) 上的作用。对任何 \(x \in L\)，用 \(\overline{x}\) 表示从 \(G\) 到 \(L/\mathfrak{m}\) 的、把 \(g\) 映到 \(\pi_{\mathfrak{m}}(gx)\) 的映射。由 VIII, 第 305 页的公式 (15) 以及 \(G_{\mathfrak{m}}\) 在 \(L/\mathfrak{m}\) 上的作用的定义，立即可得 \(\overline{x}\) 属于 \(F_{\mathfrak{m}}\)，且映射 \(u: x \mapsto \overline{x}\)（从 \(L\) 到 \(F_{\mathfrak{m}}\)）满足 \(\lambda_{F_{\mathfrak{m}}}(g) \circ u = u \circ \lambda(g)\)（对 \(g \in G\)）。换言之，\(u\) 是一个 \((K, G)\)-代数的态射。

引理 9. —— 态射 \(u\) 是满射，其核为 \(\mathfrak{a}_{\sigma}\)。由 \(u\) 通过过渡到商得到的映射是从 \(L_{\sigma}\) 到 \(F_{\mathfrak{m}}\) 的一个 (K, G)-代数同构。

由于 \(\pi_{\mathfrak{m}}\) 的核等于 \(\mathfrak{m}\)，\(u\) 的核等于 \(\bigcap_{g\in G}g^{-1}.\mathfrak{m} = \mathfrak{a}_{\sigma}\)。为证明 \(u\) 是满射，只需证明向量空间 \(\mathrm{L}_{\sigma} = \mathrm{L} / \mathfrak{a}_{\sigma}\) 与 \(\mathrm{F}_{\mathfrak{m}}\) 在 \(\mathrm{K}\) 上有相同的维数。而所有理想 \(g\cdot \mathfrak{m}\) 在 \(\mathrm{L}\) 中有相同的余维数，且由引理 8, a)，

\[
[ \mathrm{L} _ {\sigma}: \mathrm{K} ] = \operatorname{Card} (\sigma). [ \mathrm{L} / \mathfrak {m}: \mathrm{K} ].
\]

此外，由引理 7, b)，我们有

\[
[ \mathrm{F} _ {\mathfrak {m}}: \mathrm{K} ] = (\mathrm{G}: \mathrm{G} _ {\mathfrak {m}}) [ \mathrm{L} / \mathfrak {m}: \mathrm{K} ].
\]

由于我们有 \(\operatorname{Card}(\sigma)=(\mathrm{G}:\mathrm{G}_{\mathfrak{m}})\)，我们就证明了等式 \([L_{\sigma}:K]=[F_{m}:K]\)。

我们此外还假设同态 \(\lambda\)（从 G 到 \(\mathrm{Aut}_{\mathrm{K}}(\mathrm{L})\)）是单射。我们借助映射 \(\xi \mapsto \xi \cdot 1\) 把 K 与它在 L 中的像等同起来。设 \(\Omega\) 是 K 的一个代数闭扩张。集合 \(\mathcal{F}(\mathrm{G},\Omega)\)（从 G 到 \(\Omega\) 的映射组成的）与余诱导 \((\Omega ,\mathrm{G})\)-代数 \(\mathrm{Coind}_{\{e\}}^{\mathrm{G}}(\Omega)\) 重合。它是一个秩 1 的自由 \(\Omega [\mathrm{G}]\)-模。设 \(\mathcal{H}\) 是从 L 到 \(\Omega\) 的 K-代数同态的集合。我们定义 G 在 \(\mathcal{H}\) 上的一个右作用，用 \((g,\chi)\mapsto \chi \circ \lambda (g)\)。于是 \(\Omega\)-代数 \(\mathcal{F}(\mathcal{H},\Omega)\) 装备了一个 \((\Omega ,\mathrm{G})\)-代数结构（由 G 在 \(\mathcal{H}\) 上的右作用导出）。

引理 10. —— 设 L 是一个在 K 上平展的 (K, G)-代数。映射 \(\psi\)（从 \(\mathrm{L}_{(\Omega)}\) 到 \(\mathcal{F}(\mathcal{H},\Omega)\)）由关系式

\[
\psi (\xi \otimes x) = (\xi \chi (x)) _ {\chi \in \mathscr {H}}
\]

所刻画，是一个 (K, G)-代数同构。

由于 L 是平展的，映射 \(\psi\) 是一个 \(\Omega\)-代数同构（V, §6, No.~3, 第 30 页, 命题 2 与 V, §6, No.~3, 第 29 页, 命题 1, c)）。我们有关系

\[
\psi ((\operatorname{Id} \otimes \lambda (g)) (\xi \otimes x)) = (\xi (\chi \circ \lambda (g)) (x)) _ {\chi \in \mathscr {H}}
\]

对 \(\xi \in \Omega\)、\(x \in \mathrm{L}\) 与 \(g \in \mathrm{G}\)。于是 \(\psi\) 是一个 \((\Omega, \mathrm{G})\)-代数的态射。

定理 2. —— 设 G 是一个有限群，并设 L 是一个有限次交换 K-代数，它装备了由单同态 \(\lambda\)（从 G 到 \(\operatorname{Aut}_{\mathrm{K}}(\mathrm{L})\)）所给出的 G 的一个作用。下列性质等价：

\begin{enumerate}
\def\labelenumi{(\roman{enumi})}
\item
  存在 G 的一个子群 H、K 的一个有限次伽罗瓦扩张 E、从 H 到 \(\operatorname{Gal}(\operatorname{E}/\operatorname{K})\) 的一个同构，以及一个 \((\operatorname{K},\operatorname{G})\)-代数同构（从 L 到 \(\operatorname{Coind}_{\operatorname{H}}^{\operatorname{G}}(\operatorname{E})\)）。
\item
  代数 L 是平展的，且 H 是一个主齐性 G-集（I, §5, No.~6, 第 60 页, 定义 7）。
\item
  存在一个 \((\Omega,\mathrm{G})\)-代数同构 \(\psi:\mathrm{L}_{(\Omega)}\to\mathcal{F}(\mathrm{G},\Omega)\)；换言之，对每个 \(g\in G\)，自同构 \(\psi\circ\lambda(g)_{(\Omega)}\circ\psi^{-1}\)（\(\Omega^{G}\) 的）等于自同构
\end{enumerate}

\[
(x _ {h}) _ {h \in \mathrm{G}} \longmapsto (x _ {h g}) _ {h \in \mathrm{G}}.
\]

\begin{enumerate}
\def\labelenumi{(\roman{enumi})}
\setcounter{enumi}{3}
\item
  代数 \(\mathrm{L}\) 是约化的，且 \(\mathrm{L}\) 是一个秩 1 的自由 \(\mathrm{K}[\mathrm{G}]\)-模。
\item
  代数 L 是约化的，我们有 \(\operatorname{Card}(G) = [L : K]\)，且 K 是 L 的在 G 的作用下不变的元素组成的子环。
\item
  代数 L 是约化的，群 G 传递地作用在 L 的极大理想的集合上，且对 L 的每个极大理想 m，m 在 G 中的稳定子 \(G_{m}\) 忠实地作用在 L/m 上并以 K 为不变量子域。
\end{enumerate}

(i)⇒(ii)：设 E 是 K 的一个有限次伽罗瓦扩张，而 \(\tau\) 是从 H 到 \(\operatorname{Aut}_{\mathrm{K}}(\mathrm{E})\) 的一个同构。设 S 是 G 模 H 的右陪集的一个代表系。K-代数 \(F = \operatorname{Coind}_{\mathrm{H}}^{\mathrm{G}}(\mathrm{E})\) 同构于 \(\mathcal{F}(\mathrm{S}, \mathrm{E})\)（引理 7, a)）；于是它是平展的。我们用 \(\lambda_{F}\) 表示 G 在 F 上的作用。此外，设 \(\psi\) 是从 E 到 \(\Omega\) 的一个 K-代数同态，并设 \(\chi_{0}\) 是同态 \(f \mapsto \psi(f(e))\)（从 F 到 \(\Omega\)）。设 \(g \in G\) 使得 \(\chi_{0} \circ \lambda_{\mathrm{F}}(g) = \chi_{0}\)；由于 \(\psi\) 是单射，我们于是有 \(f(g) = f(e)\)（对每个 \(f \in F\)）。鉴于引理 7, a)，我们有 \(g \in H\)，且由公式

\[
f (h) = \tau (h) \cdot f (e),\tag{18}
\]

（它对每个 \(h \in H\) 成立），这只可能在 g = e 时发生。另一方面，由引理 7, b) 与 V, §10, No.~6, 第 66 页的定理 3，我们有

\[
[ \mathrm{F}: \mathrm{K} ] = (\mathrm{G}: \mathrm{H}) [ \mathrm{E}: \mathrm{K} ] = (\mathrm{G}: \mathrm{H}) \operatorname{Card} (\mathrm{H}) = \operatorname{Card} (\mathrm{G}).
\]

集合 \(\mathcal{H}\)（从 F 到 \(\Omega\) 的 K-同态组成的）的基数为 {[}F : K{]}，因为 F 是平展的（V, §6, No.~5, 第 32 页, 命题 4），所以 \(\mathrm{Card}(\mathcal{H}) = \mathrm{Card}(\mathrm{G})\)。由于由上所述 \(\chi_0\) 在 G 中的稳定子等于 \(\{e\}\)，\(\mathcal{H}\) 是一个主齐性 G-集。

(ii)⇒(iii)：假设 L 是平展的且 H 是一个主齐性 G-集。由引理 10，\((\Omega,\mathrm{G})\)-代数 \(\mathrm{L}_{(\Omega)}\) 与 \(\mathcal{F}(\mathcal{H},\Omega)\) 同构。由于 H 是一个主齐性 G-集，\((\Omega,\mathrm{G})\)-代数 \(\mathcal{F}(\mathcal{H},\Omega)\) 与 \(\mathcal{F}(\mathrm{G},\Omega)\) 同构。

(iii)⇒(iv)：假设性质 (iii) 成立。于是 \(L_{(\Omega)}\) 是代数 \(\Omega[G]\) 上的一个秩 1 自由模；后者可以典范地与 \(K[G]_{(\Omega)}\) 等同。然后我们应用 VIII, 第 37 页的定理 3。

蕴含关系 \((\mathrm{iv})\Rightarrow (\mathrm{v})\) 是直接的。

我们证明蕴含关系 \((v)\Rightarrow(vi)\)。代数 L 是约化的。由引理 8, c)，群 G 传递地作用在 L 的极大理想的集合 S 上。设 m 是 S 的一个元素。由引理 9，由于 \(\bigcap_{n\in S}n=\{0\}\)，代数 L 同构于代数 \(\operatorname{Coind}_{G_{m}}^{G}(L/m)\)。\(G_{m}\) 在 L/m 中的不变量代数于是由引理 7, c) 与 K 重合。因此同态 \(\lambda_{m}\)（从 \(G_{m}\) 到 \(\operatorname{Gal}((L/m)/K)\)）是满射。此外，由引理 7，我们有

\[
\operatorname{Card} (\mathrm{G}) = [ \mathrm{L}: \mathrm{K} ] = (\mathrm{G}: \mathrm{G} _ {\mathfrak {m}}) [ \mathrm{L} / \mathfrak {m}: \mathrm{K} ].
\]

于是 \(\operatorname{Card}(\mathrm{G}_{\mathfrak{m}}) = [\mathrm{L} / \mathfrak{m}:\mathrm{K}]\)，且同态 \(\lambda_{\mathfrak{m}}\) 是单射。

剩下要证明蕴含关系 \((\mathrm{vi})\Rightarrow(\mathrm{i})\)。设 m 是 L 的一个极大理想。由引理 9，代数 L 同构于代数 \(\operatorname{Coind}_{G_{m}}^{G}(L/m)\)（作为 \((K,G)\)-代数）。由于 \(G_{m}\) 忠实地作用在 L/m 上并以 K 为不变量子域，群同态 \(\lambda_{m}\) 定义了一个从 \(G_{m}\) 到 \(\operatorname{Gal}((L/m)/K)\) 的同构。

注记. —— 2) 在定理 2 中，我们可以把 \(\Omega\) 是代数闭的这一假设换成 \(\Omega\) 是可分闭的这一假设。事实上，若 \(L\) 是平展的，则从 \(L\) 到 \(\Omega\) 的每个 K-代数同态的像都是 \(K\) 的一个可分扩张。

\begin{enumerate}
\def\labelenumi{\arabic{enumi})}
\setcounter{enumi}{2}
\tightlist
\item
  正规基定理（V, §10, No.~9, 第 73 页, 定理 6）是定理 2 中蕴含关系 \((\mathrm{i})\Rightarrow (\mathrm{iv})\) 的一个特例。
\end{enumerate}

定义 1. —— 设 G 是一个有限群，并设 L 是 K 上的一个非零有限次交换代数，它装备了 (K,G)-代数的结构，其中 G 的作用由从 G 到群 \(\operatorname{Aut}_{\mathrm{K}}(\mathrm{L})\) 的一个单同态给出。我们说 L 是一个以 G 为群的伽罗瓦代数，如果它具有定理 2 的那些等价性质 (i) 至 (vi)。

注记. —— 4) 假设 L 是 K 的一个装备了 G 的作用 \(\lambda\) 的扩张。于是 L 是 K 上的一个伽罗瓦代数，当且仅当 L 是 K 的一个伽罗瓦扩张且 \(\lambda\) 是从 G 到 L 在 K 上的伽罗瓦群的一个同构。

\begin{enumerate}
\def\labelenumi{\arabic{enumi})}
\setcounter{enumi}{4}
\tightlist
\item
  假设群 G 是交换的。若 G 忠实且传递地作用在一个集合 X 上，则 X 的任何点的稳定子都化归为单位元（因为 X 的各点的稳定子全都相等，且它们的交化归为 G 的单位元素）。因此，定理 2 的性质 (i) 至 (vi) 也等价于下述性质：
\end{enumerate}

\begin{enumerate}
\def\labelenumi{(\roman{enumi})}
\setcounter{enumi}{6}
\tightlist
\item
  代数 \(\mathrm{L}\) 是平展的，且 \(\mathrm{G}\) 忠实且传递地作用在 \(\mathcal{H}\) 上。
\end{enumerate}

\begin{enumerate}
\def\labelenumi{\arabic{enumi})}
\setcounter{enumi}{5}
\tightlist
\item
  由 V, §10, No.~1, 第 75 页, 引理 5；V, §10, No.~1, 第 76 页, 命题 12；以及 VIII, 第 137 页, 命题 3，伽罗瓦代数是交换的且半单的。
\end{enumerate}

例. —— 1) 设 n 是一个严格正的整数，与 K 的特征指数互素。假设 K 中的 n 次单位根的群 \(\mu_{n}\) 的阶为 n.~于是对 n 的每个因子 d，d 次单位根的群 \(\mu_{d}\) 的阶为 d.~设 a 是 K 的一个非零元素，设 L 是代数 \(\mathrm{K}[\mathrm{X}]/(X^{n}-a)\)，并设 x 是 X 在 L 中的类。序列 \((1,x,\ldots,x^{n-1})\) 是体 K 上的向量空间 L 的一个基，且我们有 \(x^{n}=a\)。此外，多项式

\(X^{n}-a\) 与它的导数 \(nX^{n-1}\) 互素，于是代数 L 是平展的（V, §7, No.~2, 第 37 页, 命题 3）。对每个 \(\zeta\)（\(\mu_{n}\) 中的），自同构 \(P(X)\mapsto P(\zeta X)\)（环 K{[}X{]} 的）通过过渡到商定义了 L 的一个自同态 \(\lambda(\zeta)\)，因为我们有 \((\zeta X)^{n}-a=X^{n}-a\)；它是一个自同构。我们有

\[
\lambda (\zeta) x ^ {i} = \zeta^ {i} x ^ {i}\tag{19}
\]

对 \(0 \leqslant i < n\)。映射 \(\lambda : \zeta \mapsto \lambda(\zeta)\) 是从 \(\mu_{n}\) 到 \(\operatorname{Aut}_{\mathrm{K}}(\mathrm{L})\) 的一个单同态，且群 \(\lambda(\mu_{n})\) 在 L 中的不变量环等于 K·1。由于 \(\mu_{n}\) 的基数等于 \(n = [L : K]\)，装备了 \(\lambda\) 的作用的代数 L 是一个伽罗瓦代数（VIII, 第 308 页, 定理 2, (v)）。

设 \(r\) 是使得 \(a^r\) 属于 \(\mathrm{K}^{*n}\) 的最小严格正整数；它整除 \(n\)，且存在一个元素 \(b\)（\(\mathrm{K}^*\) 的）使得 \(a = b^{n / r}\)。于是（V, §11, No.~8, 第 91 页, 注）多项式 \(\mathrm{X}^r - b\) 不可约，且我们有 \(\mathrm{X}^n - a = \prod_{\zeta \in \mu_{n / r}} (\mathrm{X}^r - \zeta b)\)。设 \(\mathrm{E}\) 是体 \(\mathrm{K}[\mathrm{Y}] / (\mathrm{Y}^r - b)\)，并设 \(y\) 是 \(Y\) 在 \(\mathrm{E}\) 中的类。存在一个同构 \(\theta\)（从 \(\mu_{n / r}\) 到 \(\operatorname{Gal}(\mathrm{E} / \mathrm{K})\)），它由关系式 \(\theta(\xi)(y) = \xi y\) 刻画（V, §11, No.~8, 第 91 页, 例 3）。然后我们验证伽罗瓦代数 \(L\) 同构于 \((K, \mu_n)\)-代数 \(\operatorname{Coind}_{\mu_{n / r}}^{\mu_n}(E)\)。

\begin{enumerate}
\def\labelenumi{\arabic{enumi})}
\setcounter{enumi}{1}
\tightlist
\item
  现在假设体 K 的特征为 \(p \neq 0\)。设 c 是 K 的一个元素。多项式 \(f = X^{p} - X - c\) 与它的导数 \(f' = -1\) 互素，于是代数 \(L = K[X]/(f)\) 是平展的（V, §7, No.~2, 第 37 页, 命题 3）。我们用 x 表示 X 在 L 中的像；我们有关系 \(x^{p} = x + c\)。序列 \((1, x, \ldots, x^{p-1})\) 是把 L 看作 K 上的向量空间时的一个基。
\end{enumerate}

设 P 是 K 的素子域的加法群；它是一个阶为 p 的循环群，由 K 的单位元素 1 生成。对 P 中的每个 j，我们有 \(j^{p} = j\)（V, §1, No.~3, 第 4 页, 公式 (4)），因此 \(f(\mathrm{X} + j) = f(\mathrm{X})\)。故存在代数 L 的一个自同构 \(\gamma(j)\)，它由关系式 \(\gamma(j)(x) = x + j\) 刻画；此外，所得的映射 \(\gamma\) 是从 P 到 \(\operatorname{Aut}_{\mathrm{K}}(\mathrm{L})\) 的一个单同态。

设 \(\Omega\) 是 \(K\) 的一个代数闭扩张，并设 \(\xi\) 是多项式 \(f\) 在 \(\Omega\) 中的一个根。我们有 \(\xi^p = \xi + c\)，于是

\[
\mathrm{X} ^ {p} - \mathrm{X} - c = (\mathrm{X} ^ {p} - \xi^ {p}) - (\mathrm{X} - \xi) = (\mathrm{X} - \xi) ^ {p} - (\mathrm{X} - \xi) = \prod_ {j \in \mathrm{P}} (\mathrm{X} - \xi - j)
\]

由 V, §12, No.~1, 第 94 页, 公式 (1)。对 P 中的每个 \(j\)，存在唯一的代数同态 \(\chi_j: \mathrm{L} \to \Omega\)，它把 \(x\) 映到 \(\xi + j\)；此外，从 L 到 \(\Omega\) 的每个同态都是诸 \(\chi_j\) 之一，且我们有关系 \(\chi_{j}=\chi_{0}\circ\gamma(j)\)。装备了 \(\gamma\) 的代数 L 具有 VIII, 第 308 页的定理 2 的性质 (ii)，因此是 K 上的一个伽罗瓦代数。

为描述 L 的结构，我们必须区分两种情形：

\begin{enumerate}
\def\labelenumi{\alph{enumi})}
\item
  我们有 \(\xi\notin K\)。于是多项式 \(f(X)\) 在 \(K[X]\) 中不可约（V, §11, No.~9, 第 93 页, 例 3）。在这种情形，L 是 K 的一个次数为 p 的循环扩张，且 \(\gamma\) 是从 P 到 \(\operatorname{Gal}(L/K)\) 的一个同构。
\item
  我们有 \(\xi \in \mathrm{K}\)。于是映射 \(\psi : y \mapsto (\chi_j(y))_{j \in \mathrm{P}}\) 是从代数 \(\mathrm{L}\) 到积代数 \(\mathrm{K}^{\mathrm{P}}\) 的一个同构；此外，\(\psi \circ \gamma(k) \circ \psi^{-1}\) 是自同构 \((x_j)_{j \in \mathrm{P}} \mapsto (x_{j+k})_{j \in \mathrm{P}}\)（\(\mathrm{K}^{\mathrm{P}}\) 的，对每个 \(k \in \mathrm{P}\)）。
\end{enumerate}

